% ===== PRÉAMBULE CANONIQUE — à coller VERBATIM en tête de CHAQUE .tex =====
\documentclass[11pt,a4paper]{article}
\usepackage[utf8]{inputenc}\usepackage[T1]{fontenc}\usepackage{lmodern}
\usepackage[french]{babel}
\usepackage{amsmath,amssymb,mathtools}\usepackage{icomma}
\usepackage{siunitx}\sisetup{locale=FR}  % \qty{}{}, \si{}, \ang{}
\usepackage{xcolor}\usepackage{colortbl}
\usepackage{tikz}\usepackage{pgfplots}\pgfplotsset{compat=1.18}
\usepackage[siunitx,europeanresistors]{circuitikz} % circuits (élec.)
\usepackage[most]{tcolorbox}\usepackage{enumitem}\usepackage{array,booktabs}
\usepackage[a4paper,margin=2.2cm]{geometry}
\usetikzlibrary{arrows.meta,calc,angles,quotes,positioning,patterns,%
  backgrounds,shapes.geometric,decorations.pathmorphing,decorations.markings,intersections}
\definecolor{primary}{RGB}{222,49,99}\definecolor{primarydark}{RGB}{150,22,55}
\definecolor{defcol}{RGB}{0,114,178}\definecolor{methcol}{RGB}{52,92,148}
\definecolor{excol}{RGB}{0,135,90}\definecolor{attncol}{RGB}{200,40,55}
\tcbset{boxbase/.style={enhanced,breakable,boxrule=0.9pt,arc=2.5pt,
  left=3mm,right=3mm,top=2.3mm,bottom=2.3mm,fonttitle=\bfseries,coltitle=white}}
% --- boîtes TUTORIEL (noms EXACTS, ne pas renommer) ---
\newtcolorbox{cle}[1][]{boxbase,colback=defcol!6,colframe=defcol,
  title={\textsc{À retenir}\ifx\relax#1\relax\relax\else\textmd{ : \itshape #1}\fi}}
\newtcolorbox{methode}[1][]{boxbase,colback=methcol!7,colframe=methcol,sharp corners,
  title={\textsc{Méthode}\ifx\relax#1\relax\relax\else\textmd{ --- \itshape #1}\fi}}
\newtcolorbox{exemplebox}[1][]{boxbase,colback=excol!5,colframe=excol,
  title={\textsc{Exemple}\ifx\relax#1\relax\relax\else\textmd{ : \itshape #1}\fi}}
\newtcolorbox{piege}[1][]{boxbase,colback=attncol!6,colframe=attncol,
  title={\textsc{Piège}\ifx\relax#1\relax\relax\else\textmd{ : \itshape #1}\fi}}
\newtcolorbox{remarque}[1][]{boxbase,colback=black!4,colframe=black!45,
  title={\textsc{Remarque}\ifx\relax#1\relax\relax\else\textmd{ : \itshape #1}\fi}}
% --- boîtes EXERCICE / CORRIGÉ (noms EXACTS, auto-comptées) ---
\newtcolorbox[auto counter]{exo}[1][]{boxbase,colback=primary!4,colframe=primary,
  title={\textsc{Exercice}~\thetcbcounter\ifx\relax#1\relax\relax\else\textmd{ --- \itshape #1}\fi}}
\newtcolorbox[auto counter]{corrige}[1][]{boxbase,colback=excol!4,colframe=excol,
  title={\textsc{Corrigé}~\thetcbcounter\ifx\relax#1\relax\relax\else\textmd{ --- \itshape #1}\fi}}
\newcommand{\vv}[1]{\vec{#1}}\newcommand{\ds}{\displaystyle}
\newcommand{\R}{\mathbb{R}}\newcommand{\N}{\mathbb{N}}\newcommand{\Z}{\mathbb{Z}}
% (un \usepackage{fancyhdr}+en-tête est possible ; il est ignoré côté web)
% =========================================================================
\begin{document}
\sisetup{per-mode=symbol}

\begin{corrige}[travail d'une force horizontale]
\begin{enumerate}[label=\alph*)]
  \item La force est \emph{parallèle} au déplacement et dans le même sens : $\alpha=\ang{0}$, donc
        $\cos\alpha=\cos\ang{0}=1$.
  \item $W=F\times d\times\cos\alpha=200\times5\times1=\qty{1000}{\joule}$. Comme $\cos\alpha>0$, le
        travail est \textbf{moteur} ($W>0$) : la force donne de l'énergie à la caisse.
\end{enumerate}
\end{corrige}

\begin{corrige}[travail d'une force inclinée]
Seule la composante de $\vv F$ parallèle au déplacement travaille : $F_x=F\cos\ang{60}$.
\[
W=F\times d\times\cos\ang{60}=80\times10\times0,5=\qty{400}{\joule}.
\]
La force est inclinée : il ne faut \emph{pas} oublier le facteur $\cos\alpha$ (sinon on trouverait à
tort $800\ \si{\joule}$).
\end{corrige}

\begin{corrige}[une force qui ne travaille pas]
\begin{enumerate}[label=\alph*)]
  \item Le poids est \emph{vertical} et le déplacement \emph{horizontal} : ils sont perpendiculaires,
        $\alpha=\ang{90}$.
  \item $W(\vv G)=G\times d\times\cos\ang{90}=120\times15\times0=\qty{0}{\joule}$. Une force
        perpendiculaire au déplacement ne travaille pas : porter une valise en marchant à
        l'horizontale ne fournit (idéalement) aucun travail au sens mécanique.
\end{enumerate}
\end{corrige}

\begin{corrige}[travail résistant du frottement]
\begin{enumerate}[label=\alph*)]
  \item Le frottement est opposé au déplacement : $\alpha=\ang{180}$, donc $\cos\alpha=\cos\ang{180}=-1$.
  \item $W(\vv F_f)=F_f\times d\times\cos\ang{180}=30\times8\times(-1)=\qty{-240}{\joule}$. Le travail
        est \textbf{négatif} (\emph{résistant}) : le frottement \emph{retire} de l'énergie au palet et
        le freine jusqu'à l'arrêt.
\end{enumerate}
\end{corrige}

\begin{corrige}[de la puissance au travail]
\begin{enumerate}[label=\alph*)]
  \item Par définition $P=\dfrac{W}{\Delta t}$, donc $W=P\times\Delta t=750\times20=\qty{15000}{\joule}$.
  \item $W=\qty{15000}{\joule}=\qty{15}{\kilo\joule}$ (on divise par $1000$).
\end{enumerate}
\end{corrige}
\end{document}
